% rhythm: article grouped content in center: a picture the engine draws itself (one filled rectangle, the native subset, no label), a tabular, an image; each read as the whole span between its paragraphs, and the picture and image also as their edges
% boundary: par-center-picture Birch span
% boundary: par-center-picture Birch above
% boundary: center-picture-par Birch below
% boundary: par-center-tabular Cedar span
% boundary: tabular-row Dogwood span
% boundary: center-tabular-par Elm span
% boundary: par-center-image Hazel span
% boundary: par-center-image Hazel above
% boundary: center-image-par Hazel below
\documentclass{article}
\usepackage{tikz}
\usepackage{graphicx}
\usepackage{fontspec}
\setmainfont{OpenSans-Regular.ttf}[Path=../corpus/fonts/, BoldFont=OpenSans-Bold.ttf, ItalicFont=OpenSans-Italic.ttf, BoldItalicFont=OpenSans-BoldItalic.ttf]
\begin{document}
Alder is the paragraph above a centred picture.
\begin{center}
\begin{tikzpicture}
\fill (0,0) rectangle (3,1.5);
\end{tikzpicture}
\end{center}
Birch is the paragraph below the picture.
\begin{center}
\begin{tabular}{ll}
Cedar & one \\
Dogwood & two \\
\end{tabular}
\end{center}
Elm is the paragraph below the table.
\begin{center}
\includegraphics[width=2cm]{../corpus/rects.png}
\end{center}
Hazel is the paragraph below the image.
\end{document}
